% !TEX root = ../../report.tex

\section{Проверка соответствия третьей нормальной форме}

Согласно требованиям к курсовой работе, спроектированная база данных должна находиться в третьей нормальной форме (3НФ). Отношение находится в 3НФ, если оно находится во второй нормальной форме и ни один неключевой атрибут не зависит транзитивно от первичного ключа~\cite{db_models}.

Проверка выполняется последовательно для каждой нормальной формы.

\subsection*{Первая нормальная форма (1НФ)}

Отношение находится в 1НФ, если все атрибуты содержат атомарные (неделимые) значения. В спроектированной базе данных это условие выполняется:

\begin{itemize}
    \item адрес доставки разделён на атомарные компоненты: \texttt{city}, \texttt{street}, \texttt{house}, \texttt{apartment}, \texttt{zip\_code};
    \item связь <<многие ко многим>> между товарами и категориями вынесена в отдельную таблицу \texttt{product\_categories};
    \item все остальные атрибуты содержат скалярные значения (строки, числа, даты, логические значения).
\end{itemize}

\subsection*{Вторая нормальная форма (2НФ)}

Отношение находится в 2НФ, если оно находится в 1НФ и каждый неключевой атрибут полностью функционально зависит от первичного ключа (отсутствуют частичные зависимости).

В большинстве таблиц первичным ключом является суррогатный идентификатор \texttt{id} (одноатрибутный ключ), при котором частичные зависимости невозможны по определению.

Единственная таблица с составным первичным ключом --- \texttt{product\_cate\-gories}. Ключ состоит из двух внешних ключей: \texttt{product\_id} и \texttt{category\_id}. Таблица не содержит неключевых атрибутов, следовательно, условие 2НФ выполняется тривиально.

\subsection*{Третья нормальная форма (3НФ)}

Отношение находится в 3НФ, если оно находится в 2НФ и не содержит транзитивных зависимостей неключевых атрибутов от первичного ключа.

Рассмотрим таблицы, в которых потенциально могут возникнуть транзитивные зависимости:

\begin{itemize}
    \item \textbf{orders}: атрибуты \texttt{user\_id}, \texttt{seller\_id} и \texttt{address\_id} являются внешними ключами, а не производными данными. Атрибут \texttt{total\_amount} может рассматриваться как вычисляемый (сумма позиций), однако хранится денормализованно для оптимизации выборок и фиксируется на момент оформления заказа;
    \item \textbf{products}: атрибут \texttt{rating} пересчитывается триггером на основе таблицы \texttt{reviews}. Формально это производное значение, хранимое для сокращения повторных расчётов рейтинга по отзывам. Атрибут \texttt{reserved\_quantity} хранит операционное состояние резерва товара и ограничивается условием \texttt{reserved\_quantity} $\leq$ \texttt{stock\_quantity};
    \item \textbf{sellers}: аналогично, \texttt{rating} пересчитывается триггером из рейтингов товаров.
\end{itemize}

Атрибуты \texttt{rating} в таблицах \texttt{products} и \texttt{sellers}, а также \texttt{total\_amount} в таблице \texttt{orders} представляют собой контролируемую денормализацию. Их значения должны поддерживаться в согласованном состоянии правилами проектируемой базы данных и правилами обработки заказов. Все остальные неключевые атрибуты зависят только от первичного ключа своей таблицы.

Таким образом, спроектированная база данных находится в третьей нормальной форме с обоснованной денормализацией вычисляемых агрегатов.
